%%=============================================================================
%% Conclusie
%%=============================================================================

\chapter{Conclusie}%
\label{ch:conclusie}

% TODO: Trek een duidelijke conclusie, in de vorm van een antwoord op de
% onderzoeksvra(a)g(en). Wat was jouw bijdrage aan het onderzoeksdomein en
% hoe biedt dit meerwaarde aan het vakgebied/doelgroep? 
% Reflecteer kritisch over het resultaat. In Engelse teksten wordt deze sectie
% ``Discussion'' genoemd. Had je deze uitkomst verwacht? Zijn er zaken die nog
% niet duidelijk zijn?
% Heeft het onderzoek geleid tot nieuwe vragen die uitnodigen tot verder 
%onderzoek?

\lipsum[76-80]

\begin{comment}
    % Er werden tientallen werkende demo's gemaakt met verschillende technologieën waarvan maar een deel in deze corpus opgenomen ziyn
    
    % tool invocation reliability -> some questions fail to trigger a tool-use from an agent, eg.: "pls find 'x' in the documents" will get a response: "please provide the necesary documents" while it has access to a tool that can search documents, but the prompt is not a good one, it does not hint strongly enough that there is a document that the agent has a tool to use to search through, instead prompts like: "please summarize the contents of the documents" will cause the tool to be used and a summary provided, it's also possible that agents are far more familiar with certain task-phrases like summarize being linked to already existing documents, what are effective invocation phrases?
\end{comment}

